% A page contains two token positions; words label their cached K/V, not raw storage.
\begin{tikzpicture}[x=1mm,y=1mm,font=\figurefont\footnotesize,
  every node/.style={text=NNInk},
  nn flow/.append style={draw=NNWire,rounded corners=1.5mm},
  nn operator path/.append style={draw=NNWire,rounded corners=1.5mm},
  logical/.style={nn operator,minimum width=44mm,minimum height=14mm,align=center,inner sep=2pt},
  table/.style={nn transform,minimum width=32mm,minimum height=14mm,align=center,inner sep=2pt},
  physical/.style={nn frame,minimum width=38mm,minimum height=23mm,inner sep=2pt}]
  \node[nn heading,anchor=west] at (0,12) {每块容纳2个词元：逻辑顺序由块表恢复};
  \node[anchor=east] at (10,-3) {A};
  \node[logical,draw=NNInput,fill=NNInput!18] (a0) at (37,-3) {逻辑块0\\请 / 解释};
  \node[logical] (a1) at (89,-3) {逻辑块1\\注意力 / 机制};
  \node[table] (atab) at (146,-3) {A块表\\$[3,\ 1]$};
  \draw[nn flow] (a0.east)--(a1.west);
  \draw[nn operator path] (a1.east)--(atab.west);
  \node[anchor=east] at (10,-28) {B};
  \node[logical,draw=NNInput,fill=NNInput!18] (b0) at (37,-28) {逻辑块0\\请 / 解释};
  \node[logical] (b1) at (89,-28) {逻辑块1\\位置 / 编码};
  \node[table] (btab) at (146,-28) {B块表\\$[3,\ 4]$};
  \draw[nn flow] (b0.east)--(b1.west);
  \draw[nn operator path] (b1.east)--(btab.west);
  \node[nn heading,anchor=west] at (0,-57) {物理缓存：按实际地址排列};
  \foreach \i/\x in {1/20,2/62,3/104,4/146}{
    \node[physical] (p\i) at (\x,-79) {};
    \node at (\x,-72) {物理块 \i};
  }
  \node[align=center] at (20,-82) {注意力 / 机制\\A独有};
  \node[nn annotation] at (62,-82) {空闲};
  \node[draw=NNInput,fill=NNInput!18,line width=.5pt,minimum width=34mm,align=center,inner sep=2pt]
    at (104,-83) {请 / 解释\\A、B共享};
  \node[align=center] at (146,-82) {位置 / 编码\\B独有};
  \draw[NNLine,line width=.85pt] (0,-96)--(166,-96);
  \node[nn heading,anchor=west] at (0,-105) {一次具体寻址：读取A的“解释”（词元位置从0开始）};
  \node[logical,minimum width=40mm] (index) at (22,-123)
    {位置 $i=1$\\逻辑块 $\lfloor i/2\rfloor=0$};
  \node[table,minimum width=43mm] (lookup) at (83,-123)
    {查A块表索引0\\物理块编号为3};
  \node[logical,minimum width=43mm] (read) at (145,-123)
    {物理块3内偏移1\\读取“解释”的KV};
  \draw[nn flow] (index.east)--(lookup.west);
  \draw[nn flow] (lookup.east)--(read.west);
\end{tikzpicture}
